\section{Форматы данных и архитектурные подходы}

\subsection{Представление трёхмерной геометрии}
В работе используются следующие форматы:
\begin{itemize}
    \item \textbf{STEP (ISO 10303)} – стандарт для обмена данными между CAD-системами. Хранит точное параметрическое описание геометрии (B-Rep), включая грани, рёбра, вершины. Используется для точной разметки датасета и для интеграции с CAD-ядром.
    \item \textbf{STL (Stereolithography)} – триангулированное представление поверхности. Используется как вход для нейросетей благодаря своей простоте и независимости от CAD-системы.
    \item \textbf{JSON} – структурированный формат для описания операций построения. Содержит тип операции, параметры, связи с геометрическими элементами.
\end{itemize}

\subsection{Препроцессинг и нормализация}
Для обучения нейросетей STL-файлы преобразуются в облака точек или полигональные сетки с фиксированным числом рёбер. Проводится нормализация: центр модели переносится в начало координат, масштабирование к единичному размеру. Параметры операций также нормализуются (углы в диапазон 0–1, длины масштабируются).

\subsection{Архитектуры для анализа 3D-геометрии}
В литературе предложены различные нейросетевые архитектуры для работы с трёхмерными данными:
\begin{itemize}
    \item \textbf{PointNet / PointNet++} – работают непосредственно с облаками точек, обеспечивая инвариантность к перестановке точек.
    \item \textbf{DGCNN (Dynamic Graph CNN)} – строит динамический граф в пространстве признаков, эффективно извлекая локальную структуру.
    \item \textbf{MeshCNN} – оперирует рёбрами треугольной сетки, используя свёртки и пулинг, сохраняющие топологию. Эта архитектура легла в основу разработанного решения благодаря её способности работать с полигональными сетками и выделять признаки на рёбрах.
\end{itemize}

